% New mathematical fragment for synthesized naming AT OPENING.
% The enclosing report supplies standard theorem/lemma/corollary environments.
\section{Synthesizing a witness name when its clause is opened}
\label{sec:named-opening}

The new rule substitutes its inferred numerical name into
the opened clause, so its branch proves the desired outer conclusion
directly. The rule requires neither an equality proof, a supplied
replacement term, nor a separate branch-result formula. The main
completeness issue is admissibility of substitution. We resolve it before
recovering unrestricted existential-style elimination, using only two
guard-free forms of elimination that are immediately available.

\subsection{Complete presentation}

Numerical terms are full computation normal forms of
\[
 t::=x\mid0\mid S(t)\mid I(n,a;u.h)
\]
under
\[
 I(0,a;u.h)\longrightarrow a,\qquad
 I(Sn,a;u.h)\longrightarrow h[I(n,a;u.h)/u].
\]
Only $u$ in $h$ is bound. Write $N$ for normalization. Substitution is
simultaneous, capture avoiding, and followed by $N$. Inserting a term
does not recursively substitute into that inserted copy. The inherited
termination, confluence, and substitution-coherence results for $N$ are
part of the metatheoretic foundation and of the trusted checker boundary.

Formulas and optional-premise clauses are
\[
 A::=(s=t)\mid D_{x,y}B,\qquad
 [\Delta;y.B(y)],\quad |\Delta|\in\{0,1\},
 \quad y\notin\operatorname{FV}(\Delta).
\]
Each context is a finite list of clause occurrences containing the fixed
axiom occurrence $\mathcal R=[\,;y.y=y]$. An ordinary assumption $A$
means the unconditional constant clause $[\,;v.A]$. Put $H=(0=0)$,
$1=S0$, and $2=S(S0)$.

The intended ordinary readings are actual numerical equality,
\[
 (D_{x,y}B)^\circ=\exists x\bigl(B^\circ(x,0)\to
                                      \forall y\,B^\circ(x,y)\bigr),
\]
and $[\,;y.B]^\circ=\forall y B^\circ(y)$,
$[[A];y.B]^\circ=A^\circ\to\forall y B^\circ(y)$.
The three primitive inference families are the following.

\paragraph{Use.}
For a retained clause $[\Delta;y.B]\in\Gamma$,
\[
 \frac{\{\,\Gamma\vdash A:A\text{ is the antecedent occurrence in }
                                  \Delta\,\}}
      {\Gamma\vdash N(B(t))}.
 \tag{U}
\]
An empty antecedent requires no proof; a singleton requires one.
In particular, a use of $\mathcal R$ proves $t=t$.

\paragraph{Introduce.}
\[
 \frac{\Gamma,B(t,\eta)\vdash B(t,Sz)}
      {\Gamma\vdash D_{x,y}B(x,y)},\qquad
 \eta=\begin{cases}
 z,&B\text{ has equation root and }t\equiv0,\\
 0,&\text{otherwise}.
 \end{cases}
 \tag{I}
\]
The witness $t$ is canonical and $z$ is fresh for the retained context,
the witness, and all passive matrix parameters. Literal syntax, not
provable equality, determines the mode.

\paragraph{Open with a synthesized name.}
For a major $D_{x,y}B(x,y)$ choose a fresh parameter $b$ and put
\[
 B_b(y)=N(B(b,y)),\qquad
 \tau_B(b)=\begin{cases}
 r,&B_b(0)\equiv H\text{ and }B_b(1)\equiv(b=r),\\
 b,&\text{otherwise}.
 \end{cases}
 \tag{P}
\]
Define the substituted clause
\[
 \mathcal C^*_B(b)=
 [[N(B(\tau_B(b),0))];y.N(B(\tau_B(b),y))].
\]
Then the rule is
\[
 \frac{\Gamma\vdash D_{x,y}B(x,y)\qquad
       \Gamma,\mathcal C^*_B(b)\vdash A}
      {\Gamma\vdash A},
 \qquad b\notin\operatorname{FV}(\Gamma,A).
 \tag{O$^*$}
\]
The parameter $b$ is also fresh for the passive parameters of $B$.
The synthesized $r$ may contain $b$; the branch still lives in the scope
with $b$ available, and its conclusion must be in the outer scope. The
primitive rule has exactly two proof premises. There is no additional
computation, equality, hypothesis, induction, cut, or substitution rule.
Computation is part of canonical syntax, and reflexivity is the explicit
fixed axiom clause.

\begin{lemma}[Direct ordinary soundness]
\label{lem:opening-sound}
O$^*$ is sound in intuitionistic first-order equality logic over the scalar
computation equations.
\end{lemma}
\begin{proof}
Open the existential major at $b$, obtaining
$\mathcal C_B(b)=[[B_b(0)];y.B_b(y)]$. If P selects its first branch,
the zero instance is reflexive and clause use at $1$ gives $b=\tau_B(b)$.
Otherwise reflexivity gives the same equality because $\tau_B(b)=b$.
Ordinary equality transport converts the whole clause reading
$\mathcal C_B(b)^\circ$ to
$\mathcal C^*_B(b)^\circ$. Apply the supplied branch derivation to this
converted assumption. The conclusion contains no $b$, so ordinary
existential elimination closes. Equality transport and simultaneous
substitution are valid also for an equality $b=r(b)$.
\end{proof}

\subsection{The preliminary rules that do not need general substitution}

Weakening by clause occurrences, injective fresh renaming, and cut of an
ordinary formula assumption are admissible by direct induction on
derivations. For renaming, the injectivity and freshness preserve and
reflect the literal patterns in I and P; normalization is equivariant.
For ordinary-assumption cut, replace only uses of the selected
unconditional constant clause, retain all other clause occurrences, and
freshen local parameters before insertion. These arguments do not
substitute arbitrary numerical terms into a derivation.

\begin{lemma}[Two immediately ordinary forms of opening]
\label{lem:opening-safe}
O$^*$ is exactly ordinary paired elimination in each of the following cases:
\begin{enumerate}
 \item the witness coordinate $x$ is unused in the matrix $B(x,y)$;
 \item the matrix $B$ has compound, rather than equation, root.
\end{enumerate}
\end{lemma}
\begin{proof}
In the first case $B_b(1)$ contains no $b$, so its left-hand term cannot
literally be the fresh parameter $b$. In the second case $B_b(0)$ still
has compound root, because term normalization does not erase formula
binders, and therefore cannot be $H$. In either case P returns $b$.
\end{proof}

For unused $w$ write $q_yA(y)=D_{w,y}A(y)$.
\begin{lemma}[Closures and numerical cases]
\label{lem:opening-cases}
Before using general ordinary elimination or equality transport, the
following rules are available:
\[
 \frac{\Gamma,A(0)\vdash A(Sz)}{\Gamma\vdash q_yA(y)},
 \qquad
 \frac{\Gamma\vdash q_yA(y)\qquad\Gamma\vdash A(0)}
      {\Gamma\vdash A(t)}.
\]
Consequently $\Gamma\vdash A(0)$ and $\Gamma\vdash A(Sz)$, with $z$
fresh, imply $\Gamma\vdash A(t)$ for every well-scoped $t$.
\end{lemma}
\begin{proof}
Introduce the unused witness of $q$ at $1$, forcing I's case mode.
For specialization, use the unused-witness case of
Lemma~\ref{lem:opening-safe}, which opens the clause
$[[A(0)];y.A(y)]$; U at $t$ with the given base proof gives the desired
outer conclusion. The numerical-case rule follows by weakening the
successor proof with the unused assumption $A(0)$ and composing the
two constructions.
\end{proof}

\begin{lemma}[Unrestricted case introduction]
\label{lem:opening-case-repair}
For every canonical matrix $B$ and every canonical witness $t$, the rule
\[
 \frac{\Gamma,B(t,0)\vdash B(t,Sz)}
      {\Gamma\vdash D_{x,y}B(x,y)}
\]
is derivable, with the same freshness condition on $z$ as in I. Its
derivation uses only U, I, and the two ordinary opening forms of
Lemma~\ref{lem:opening-safe}.
\end{lemma}
\begin{proof}
The primitive I already suffices unless $B$ has equation root and $t=0$.
Write the remaining matrix as $E(x,y)$. Put
\[
 C(v,a,c)=N(I(v,a;u.c)),\quad u\notin\operatorname{FV}(c),
\]
and for an equation $E=(s=t)$ define
\[
 G_vE=(C(v,0,s)=C(v,0,t)),\qquad L(E)=q_vG_vE.
\]
The literal endpoints are $G_0E=H$ and $G_{Sw}E=E$.
Lemma~\ref{lem:opening-cases} gives $E\vdash L(E)$ by closure
introduction and $L(E)\vdash E$ by specialization at $1$ and reflexivity
of $H$. Only unused-witness opening is involved.

Let $A(x,y)=L(E(x,y))$. Under $A(0,0)$, unwrap to $E(0,0)$, insert the
supplied premise by ordinary-assumption cut, and wrap $E(0,Sz)$ as
$A(0,Sz)$. Since $A$ is compound, primitive I at zero uses cases and
gives $D_{x,y}A(x,y)$. Open this major at fresh $c$ using the
compound-matrix case of Lemma~\ref{lem:opening-safe}. To introduce
$D_{x,y}E(x,y)$ at witness $c$, assume $E(c,0)$, wrap it to $A(c,0)$,
use the opened clause at $Sz$, and unwrap $A(c,Sz)$ to $E(c,Sz)$.
The fresh witness $c$ is not literal zero, so I uses cases. The resulting
$D_{x,y}E(x,y)$ is independent of $c$ and can be exported. No general
substitution theorem and no unrestricted opening rule was used.
\end{proof}

\subsection{The substitution theorem and its noncircular proof}

A simultaneous numerical substitution may send a parameter to a term
containing that same parameter. Images are inserted once. Only the finitely
many parameters appearing in a given derivation matter. As usual, bound
variables and eigenparameters are freshened away from the images before
the corresponding substitution is carried out.

\begin{theorem}[Admissibility of arbitrary numerical substitution]
\label{thm:opening-substitution}
If $\Gamma\vdash A$ and $\sigma$ is any well-scoped simultaneous
numerical substitution, then
\[
 N(\Gamma\sigma)\vdash N(A\sigma).
\]
This is an effective transformation of finite derivations. It is not a
claim that substitution preserves every primitive rule instance.
\end{theorem}
\begin{proof}
Use strong induction on the original derivation, with the induction
assertion quantified over \emph{all} simultaneous substitutions. This
quantification is essential: in one opening case the branch receives a
different combined substitution from the major. The induction hypothesis
is applied to original proper subderivations, never to an already
translated derivation of potentially larger size.

For U, substitute its selected clause occurrence and its zero or one
premises, then use normalization--substitution coherence. The distinguished
axiom $\mathcal R$ is unchanged because it has no free numerical parameter.
For I, first freshen $z$ away from all images of $\sigma$. A progress
instance stays progressive: its equation root and literal-zero witness
are preserved. A case instance stays a case instance unless an equation
matrix has a previously nonzero witness whose substituted normal form
is zero. In that exceptional case, the induction hypothesis gives exactly
the premise of the derived unrestricted case rule in
Lemma~\ref{lem:opening-case-repair}; apply that rule. That lemma was proved
before the present substitution theorem.

It remains to treat O$^*$. Freshen its parameter $b$ away from the passive
data and from the finitely many substitution images. Write
$B'=N(B\sigma)$ and $\tau'=\tau_{B'}(b)$. Translate the major by the
induction hypothesis. There are two cases according to the \emph{old}
probe.

If the old probe selected $r$, its two literal identities were
$B_b(0)=H$ and $B_b(1)=(b=r)$. Applying $\sigma$ while leaving $b$ fixed
preserves these identities after normalization: $H$ is closed and $b$
is untouched. Formula roots are preserved. Therefore the new probe is
also successful and $\tau'=N(r\sigma)$ literally. Apply the induction
hypothesis to the original branch using $\sigma$ with $b$ fixed. Its
substituted clause is, by coherence,
\[
 [[B'(\tau',0)];y.B'(\tau',y)],
\]
which is precisely the clause needed by the new O$^*$ instance.

If the old probe used its fallback, its branch clause was the ordinary
$[[B(b,0)];y.B(b,y)]$. The new probe may remain false or may become true.
In \emph{both} cases apply the induction hypothesis once to the original
branch with the combined simultaneous substitution that agrees with
$\sigma$ on the outer parameters and sends $b$ to $\tau'$. Its image
$\tau'$ is inserted once, including when it contains $b$. The resulting
clause is exactly
$[[B'(\tau',0)];y.B'(\tau',y)]$. The outer context and branch conclusion
become $N(\Gamma\sigma)$ and $N(A\sigma)$ because neither originally
contained $b$, and the freshening ensured that no outer image introduces
$b$. This supplies the new branch whether or not the guard was newly
created. Close the new O$^*$ instance.

This exhausts the rules. The only proof transformation invoked outside
the induction hypotheses was the already established case repair.
\end{proof}

\begin{corollary}[Unrestricted ordinary paired elimination]
\label{cor:opening-ordinary}
For arbitrary canonical $B$, the rule
\[
 \frac{\Gamma\vdash D_{x,y}B(x,y)\qquad
       \Gamma,[[B(b,0)];y.B(b,y)]\vdash A}
      {\Gamma\vdash A},\qquad
 b\notin\operatorname{FV}(\Gamma,A),
\]
is admissible, where $b$ is also fresh for the passive parameters of $B$.
\end{corollary}
\begin{proof}
Apply Theorem~\ref{thm:opening-substitution} to the supplied branch,
with the simultaneous substitution $b\mapsto\tau_B(b)$ and all outer
parameters fixed. Since $\Gamma,A$ are $b$-free, the result is exactly
the branch of O$^*$ with the same outer conclusion $A$.
The major is unchanged. This proves general ordinary elimination only
\emph{after} the substitution theorem, avoiding circularity.
\end{proof}

\subsection{Recovering every equality rule}

With $C$ as above define
\[
 P(y)=N(I(y,0;u.C(u,1,0))).
\]
Computation gives $P(1)=1$ and $P(2)=0$ literally; no parity induction is
needed.

\begin{theorem}[Uncrossed transport]
\label{thm:opening-transport}
For every fixed numerical template $h(w)$,
\[
 \Gamma\vdash s=t,\quad\Gamma\vdash N(h(s))=r
 \quad\Longrightarrow\quad\Gamma\vdash N(h(t))=r.
\]
\end{theorem}
\begin{proof}
Use the canonical equation matrix
\[
 B(x,y)=
 \bigl(C(y,0,C(P(y),N(h(x)),x))
      =C(y,0,C(P(y),r,t))\bigr).
\]
Its zero, first, and second rows are respectively
\[
 H,\qquad x=t,\qquad N(h(x))=r.
\]
At witness $s$, numerical cases for
$C(v,N(h(s)),s)=C(v,r,t)$, with the two supplied premises as its zero
and successor cases, give the successor row at $v=P(Sz)$. This ignores
I's local assumption, so I introduces $D_{x,y}B(x,y)$ in either mode.
Open that major. Its probe returns $t$, so its branch clause is
\[
 [[H];y.B(t,y)].
\]
Reflexivity proves $H$, and U at $2$ proves $N(h(t))=r$ directly.
That equation is in the outer scope, so O$^*$ closes. The numerical-case
construction was derived before equality, and the transport construction
does not require a general substitution of the supplied equality proofs.
\end{proof}

Symmetry follows by taking $h(w)=w$ and the reflexive minor $s=s$;
crossed transport follows by symmetry of the output. Transitivity and
term congruence follow in the usual finite equational constructions.
Replacing both sides of an equation uses two one-term transports and
symmetry.

Full formula transport follows by external induction on formula depth.
For a compound template $D_{x,y}B(v,x,y)$, open the source formula by
Corollary~\ref{cor:opening-ordinary} at a fresh $c$, and introduce the
target formula at witness $c$, which forces case mode. Under its zero
assumption, transport backward in the smaller matrix, specialize the
source clause at $Sz$, and transport forward in the smaller matrix.
Close introduction and ordinary elimination. This uses no object-theory
induction on quantified formulas.

\subsection{Identity-language equivalence and both interpretation composites}

Let $\mathcal T$ be the inherited fully normalized clause calculus with
the same syntax, Ax, U and I, together with ordinary paired elimination
and crossed one-term transport. Let $\mathcal O$ be the new three-rule
calculus U, I, O$^*$.

\begin{theorem}[Exact equivalence on arbitrary clause sequents]
\label{thm:opening-equivalence}
For every canonical finite clause context $\Gamma$ containing Ax and
every canonical formula $A$,
\[
 \Gamma\vdash_{\mathcal O}A
 \quad\Longleftrightarrow\quad
 \Gamma\vdash_{\mathcal T}A.
\]
Both translations are effective and leave the full end sequent literally
unchanged.
\end{theorem}
\begin{proof}
From $\mathcal T$ to $\mathcal O$, translate proof trees recursively,
replacing ordinary elimination by Corollary~\ref{cor:opening-ordinary}
and crossed transport by Theorem~\ref{thm:opening-transport} and derived
symmetry. U and I are shared. The preliminary admissibility theorems were
proved internally in $\mathcal O$ and do not rely on this translation.

Conversely, $\mathcal T$ has full formula transport by the same
formula-depth argument. Translate an O$^*$ major and open it ordinarily
at $b$, retaining its clause $\mathcal C_B(b)$. Derive
$b=\tau_B(b)$ by reflexivity and U at $1$ when the probe succeeds, or
by reflexivity alone otherwise. In this opened context derive
$q_y B(\tau_B(b),y)$ as follows. Under its zero assumption
$B(\tau_B(b),0)$, transport backward along the derived equality to
$B(b,0)$; use the old clause at $Sz$; transport forward to
$B(\tau_B(b),Sz)$; and introduce the unused $q$ witness at $1$.
Ordinary elimination of that $q$ opens exactly
$\mathcal C^*_B(b)$. Insert the translated O$^*$ branch, weakened by
the still-retained old clause, and close this unused-witness elimination.
The branch conclusion is free of $b$, so the outer ordinary elimination
also closes. All operations are on finite proof trees, with fresh
parameters chosen before insertion.
\end{proof}

\begin{corollary}[Full iPRA and the identity numerical composites]
\label{cor:opening-ipra}
The new calculus is equiconsistent and bi-interpretable with full
first-order intuitionistic primitive recursive arithmetic. Both
interpretations use the whole numerical domains and actual equality;
both numerical composite isomorphisms have graph $x=y$.
\end{corollary}
\begin{proof}
Compose Theorem~\ref{thm:opening-equivalence} with the inherited canonical
clause/native-scalar comparison and the inherited scalar/iPRA
bi-interpretation. The new comparison is identity on every numerical term
and formula, so it changes neither numerical map nor the formula
translations. Each new composite therefore consists of identity legs
and an inherited composite already provably equivalent to the original
formula, with numerical graph $x=y$. The designated contradiction
$0=S0$ is unchanged, giving equiconsistency as a special case.
\end{proof}

The new universal metatheorems are
Theorems~\ref{thm:opening-substitution} and
\ref{thm:opening-equivalence}. Their dependency order is: structural
renaming and ordinary-assumption cut; the two guard-free openings;
closure rules; case repair; strong-induction substitution; unrestricted
ordinary opening; equality; exact comparison. The iPRA arithmetic
compiler and normalization foundation remain explicitly inherited
results. No newly proof-assistant-checked status or primitive-recursive
bound on full normalization is asserted.
